\paragraph{Linear contours that mix the indices, and real spacelike momenta (proved at \(N=2\)).}
Keep Definition~(A), \(\gamma=0\) and the notation of the two preceding paragraphs, and take the contours \(A_\mu=\sum_\nu\Lambda_{\mu\nu}X_\nu\) with \(\Lambda\in\mathrm{GL}(10,\mathbb C)\) and \(X_\nu\) traceless Hermitian. They contain the \(\mathcal R_u\) and their images under complex Lorentz transformations. Put \(h=\Lambda^{\mathrm T}\eta\Lambda\) and write \(h(a,b)=a^{\mathrm T}hb\). At \(N=2\), with \(X_\nu=\sum_aX_\nu^a\sigma^a/2\) and \(x^a=(X_\nu^a)_\nu\in\mathbb R^{10}\), the matrix \(\mathcal G\) of the preceding paragraphs is \(\mathcal G_{ab}=h(x^a,x^b)\), and \(S_{\mathrm b}=\tfrac1{4g_{\mathrm{YM}}^2}\sum_{a<b}[h(x^a,x^a)h(x^b,x^b)-h(x^a,x^b)^2]\), which for \(h=\delta\) is \(S_{\mathrm E}\). Hence \(\operatorname{Re}\mathrm iS_{\mathrm b}\le-c\,S_{\mathrm E}(X)\) holds on the contour for some \(c>0\) exactly when \(h\) lies in the set \(\mathfrak C\) defined by
\[
\operatorname{Im}\bigl[h(a,a)h(b,b)-h(a,b)^2\bigr]>0\quad\text{for all linearly independent real }a,b,
\]
and then, by the Cauchy--Binet formula, the integrand at \(\epsilon=0\) times a bounded \(O\) is dominated as in the bound along the flat directions. For \(N\ge3\) the same inequality, applied to \(X\) in an \(\mathfrak{su}(2)\) block, implies \(h\in\mathfrak C\).

(i) \emph{Normal form.} Let \(h\in\mathfrak C\). Then \(h(a,a)\neq0\) for real \(a\neq0\), since otherwise a real \(b\notin\mathbb Ra\) with \(\operatorname{Re}h(a,b)=0\) exists, and for it the bracket is \(-h(a,b)^2\), which is real. As the unit sphere of \(\mathbb R^{10}\) is simply connected, \(\arg h(a,a)\) has a continuous branch on it, which is even. On the unit circle of each real 2-plane the values \(h(a,a)\) form an ellipse that avoids \(0\) and around which this branch winds zero times, so the ellipse lies in an open half-plane and the branch varies by less than \(\pi\) on the circle. The plane through a maximum and a minimum then shows that it varies by less than \(\pi\) on the sphere. So \(\operatorname{Re}(\mathrm e^{-\mathrm i\theta}h)\) is positive definite for some \(\theta\), and a simultaneous real diagonalisation gives \(h=T^{\mathrm T}\operatorname{diag}(d_\mu)T\) with \(T\) real and invertible. The condition defining \(\mathfrak C\) becomes \(\operatorname{Im}(d_\mu d_\nu)>0\) for \(\mu\neq\nu\), and with the arguments \(\beta_\mu\) of the \(d_\mu\) taken in \((\theta-\pi/2,\theta+\pi/2)\), all sums \(\beta_\mu+\beta_\nu\) lie in one interval \((2k\pi,(2k+1)\pi)\). So \(\mathfrak C=\mathfrak C_0\cup(-\mathfrak C_0)\), where \(\mathfrak C_0\) consists of the \(T^{\mathrm T}\operatorname{diag}(d_\mu)T\) whose arguments can be chosen with \(0<\beta_\mu+\beta_\nu<\pi\) for \(\mu\neq\nu\).

(ii) \emph{The class adjacent to the Hermitian contour.} On \(\mathfrak C_0\) at most one \(\beta_\mu\) lies outside \((-\pi/2,\pi/2)\) and at most one outside \((0,\pi)\). So \(\operatorname{Re}h\) and \(\operatorname{Im}h\) each have at least nine positive eigenvalues there, and at least nine negative ones on \(-\mathfrak C_0\). At a common limit point both would have kernels of dimension at least eight, so \(h\) would be singular. Since \(h=\Lambda^{\mathrm T}\eta\Lambda\) is invertible, a path of these contours whose \(h\) stays in the closure of \(\mathfrak C\) cannot pass from \(\mathfrak C_0\) to \(-\mathfrak C_0\). The Hermitian contour, \(h=\eta\), is the limit as \(u\to0\) of the \(\mathcal R_u\), which lie in \(\mathfrak C_0\) for \(0<u<2\). It is not a limit of points of \(-\mathfrak C_0\), since there \(\operatorname{Re}h\) has at least nine negative eigenvalues.

(iii) \emph{Momenta reached.} For \(k\in\mathbb C^{10}\), the observable \(\operatorname{Tr}\mathrm e^{\mathrm ik^\mu A_\mu}\) is bounded on the contour exactly when \(q=\Lambda^{\mathrm T}k\) is real, and then \(k^\mu k_\mu=q^{\mathrm T}h^{-1}q=\sum_\mu(T^{-\mathrm T}q)_\mu^2/d_\mu\). For \(h\in\mathfrak C_0\) the pairwise conditions, applied with a third index, put the arguments \(-\beta_\mu\) of the directions \(\mathrm e^{-\mathrm i\beta_\mu}\) in an interval of length less than \(\pi\) inside \((-\pi,\pi/2)\), so \(k^\mu k_\mu\neq0\) lies in the sector \(-\pi<\arg(k^\mu k_\mu)<\pi/2\). This is the sector \(|\arg t|<3\pi/4\) for \(t=\mathrm e^{\mathrm i\pi/4}k^\mu k_\mu\), and the \(\mathcal R_u\) already reach every point of it. No contour of this class with \(h\in\mathfrak C_0\) therefore makes the Fourier mode bounded for \(k^\mu k_\mu\) outside the sector, in particular for real \(k\neq0\) with \(k^\mu k_\mu\le0\). By the inclusion above, this holds for every \(N\).

(iv) \emph{Values at \(N=2\).} For \(h\in\mathfrak C\) and real \(q\), let \(W(h,q)\) be the integral over real \(X\) of \((\det\mathcal G)^4\mathrm e^{\mathrm iS_{\mathrm b}}\operatorname{Tr}\mathrm e^{\mathrm iq^\nu X_\nu}\), and let \(V(h,q)=2W(h,q)/W(h,0)\). Then \(W\) is holomorphic in \(h\), is even in \(h\), and satisfies \(W(h,q)=|\det T|^3W(T^{\mathrm T}hT,T^{\mathrm T}q)\) for real invertible \(T\). On \(\mathfrak C_0\) this gives \(W(T^{\mathrm T}\operatorname{diag}(d_\mu)T,0)=c\,|\det T|^{-3}\prod_\mu d_\mu^{-3/2}\) for a continuous branch, with \(c\neq0\), and \(V(T^{\mathrm T}\operatorname{diag}(d_\mu)T,q)=U(w)\) with \(w_\mu=(T^{-\mathrm T}q)_\mu^2/d_\mu\), where \(U\) is holomorphic on the connected set of \(w\) for which the arguments of the \(1/w_\mu\) can be chosen with pairwise sums in \((0,\pi)\); zero components of \(T^{-\mathrm T}q\) follow by continuity in \(q\). For \(w=\mathrm e^{-\mathrm i\pi/4}x\) with \(x\in(0,\infty)^{10}\) the contour is, after a real rescaling of the \(X_\nu\), the Euclidean one with \(q_\mu=\sqrt{x_\mu}\), and \(\mathrm{SO}(10)\) invariance gives \(U(w)=F_{\mathrm E}(\sum_\mu x_\mu)\). That set of \(w\) is open in a maximal totally real subspace, \(\sum_\mu w_\mu\) lies in the sector of (iii), and \(F_{\mathrm E}\) is holomorphic on \(|\arg t|<3\pi/4\) by the preceding paragraph, so the identity theorem gives
\[
V(h,q)=F_{\mathrm E}\bigl(\mathrm e^{\mathrm i\pi/4}q^{\mathrm T}h^{-1}q\bigr)\quad(h\in\mathfrak C_0),\qquad V(h,q)=F_{\mathrm E}\bigl(-\mathrm e^{\mathrm i\pi/4}q^{\mathrm T}h^{-1}q\bigr)\quad(h\in-\mathfrak C_0).
\]

(v) \emph{Pointwise limits at \(N=2\).} Let \(\Lambda(\tau)\), \(0\le\tau\le1\), be piecewise continuously differentiable, with \(\Lambda(0)\) real, \(h(\tau)\) in the closure of \(\mathfrak C\), \(\operatorname{Re}\Lambda(\tau)^{\mathrm T}\Lambda(\tau)\) positive definite, \(h(1)\in\mathfrak C_0\) and \(\Lambda(1)^{\mathrm T}k\) real. Stokes' theorem at fixed \(\epsilon\), as for the \(\mathcal R_u\), and dominated convergence on the last contour then give \(\lim_{\epsilon\to0^+}\langle\operatorname{Tr}\mathrm e^{\mathrm ik^\mu A_\mu}\rangle_\epsilon=F_{\mathrm E}(\mathrm e^{\mathrm i\pi/4}k^\mu k_\mu)\). For real \(k\) with \(k^\mu k_\mu=\kappa^2>0\) such a path exists. Take a real Lorentz transformation \(L\) with \(L^{\mathrm T}k=(0,\kappa,0,\dots,0)\), phases \(\beta(\tau)=(\pi-3\pi u\tau/4,\,-\pi u\tau/8,\,\pi u\tau/4,\dots,\pi u\tau/4)\), which give \(h(\tau)\) in the closure of \(\mathfrak C_0\) and \(h(1)\in\mathfrak C_0\) for \(0<u<8/7\), and \(\Lambda(\tau)=LS(\tau)\operatorname{diag}(\mathrm e^{\mathrm i(\beta_0-\pi)/2},\mathrm e^{\mathrm i\beta_1/2},\dots,\mathrm e^{\mathrm i\beta_9/2})\). Here \(r\) is real in the \((1,2)\)-plane with \(r_1^2\mathrm e^{-\mathrm i\beta_1}+r_2^2\mathrm e^{-\mathrm i\beta_2}=1\), and \(S(\tau)\) is a complex Lorentz transformation, obtained by Gram--Schmidt for the bilinear form \(\eta\), whose row \(\mu=1\) has the entries \(r_\nu\) divided by the \(\nu\)-th diagonal entry; the condition on \(r\) gives this row \(\eta\)-norm \(1\). Then \(h(\tau)=\operatorname{diag}(\mathrm e^{\mathrm i\beta_\mu(\tau)})\) and \(\Lambda(\tau)^{\mathrm T}k=\kappa r\). As \(u\to0\) the path tends uniformly to a real matrix, so \(\operatorname{Re}\Lambda^{\mathrm T}\Lambda\) stays positive definite for small \(u\). Hence for every real spacelike \(k\), in any direction, the pointwise limit exists and equals \(F_{\mathrm E}(\mathrm e^{\mathrm i\pi/4}k^\mu k_\mu)\), a value on the ray \(\arg t=\pi/4\) inside the sector. It depends only on \(k^\mu k_\mu\), although the regulator breaks the boosts, and the same holds with any positive definite \(M\) in place of the regulator.

For linear contours, this order of limits and \(\gamma=0\), the contour integrals with \(h\in\mathfrak C_0\) and bounded Fourier modes are therefore at \(N=2\) values of \(F_{\mathrm E}\) at \(\mathrm e^{\mathrm i\pi/4}k^\mu k_\mu\) inside the sector, they are the pointwise limits of the regulated integral wherever a path as in (v) exists, which includes every real spacelike \(k\), and the real momenta that no contour with \(h\in\mathfrak C_0\) reaches are exactly the nonzero causal ones, \(k^\mu k_\mu\le0\). Contours with \(h\in-\mathfrak C_0\) do make the mode bounded at real timelike \(k\), for example \(A_0=X_0\), \(A_i=\mathrm e^{2\pi\mathrm i/3}X_i\) at \(k^\mu=\kappa\delta^\mu_0\), with the value \(F_{\mathrm E}(\mathrm e^{\mathrm i\pi/4}\kappa^2)\) at \(N=2\), but by (ii) no path along which \(h\) stays in the closure of \(\mathfrak C\) joins them to the Hermitian contour. This is no statement about non-linear contours, about whether \(F_{\mathrm E}\) continues beyond \(\arg t=-3\pi/4\), about the pointwise limit at real timelike or null \(k\), about paths along which \(\operatorname{Re}\mathrm iS_{\mathrm b}\le0\) but \(h\) leaves the closure of \(\mathfrak C\), about (iv) and (v) for \(N\ge3\), about \(N\to\infty\) taken first, or about \(\gamma\neq0\) and Definitions~(B) and~(C). The light cone found here is the one of \(\eta\) in \(S_{\mathrm b}\), so no signature is derived and Assumption~\ref{ass:lorentzian} is not discharged.
